% ============================================================
%  ANSWER KEY — Practice 6: Combinations
% ============================================================

\section*{Վարժություն 6. Զուգորդություններ և Նյուտոնի երկանդամ}

\subsubsection*{Խնդիր 1}
9 երկարությամբ բիթային տողեր՝ ճիշտ 5 հատ զրոներով (և 4 հատ մեկերով):

Ընտրել դիրքեր 5 զրոների համար՝ $\binom{9}{5}$:

$\boxed{\binom{9}{5} = 126}$

\subsubsection*{Խնդիր 2}
Կոմիտե՝ կազմված 6 ուսուցչից և 8 աշակերտից, ընտրված 12 ուսուցիչներից և 20 աշակերտներից:

$\boxed{\binom{12}{6}\cdot\binom{20}{8} = 924 \times 125970 = 116{,}396{,}280}$

\subsubsection*{Խնդիր 3}
Ընտրել 13 խաղաքարտ 52-ից այնպես, որ դրանք պարունակեն նույն մասի (suit) 5 խաղաքարտ:

Օգտագործելով լրացումը (ոչ մի մաս չունի $\ge 5$ խաղաքարտ, այսինքն՝ յուրաքանչյուր մաս տալիս է $\le 4$ խաղաքարտ):

Հնարավոր բաշխումները $(a,b,c,d)$, որտեղ յուրաքանչյուրը $\le 4$ է և գումարը 13 է՝ $(4,4,4,1)$, $(4,4,3,2)$, $(4,3,3,3)$ են:

\begin{align*}
(4,4,4,1):&\quad 4\cdot\binom{13}{4}^3\binom{13}{1} = 19{,}007{,}345{,}500\\
(4,4,3,2):&\quad 12\cdot\binom{13}{4}^2\binom{13}{3}\binom{13}{2} = 136{,}852{,}887{,}600\\
(4,3,3,3):&\quad 4\cdot\binom{13}{4}\binom{13}{3}^3 = 66{,}905{,}856{,}160
\end{align*}

$N_{\le 4} = 222{,}766{,}089{,}260$:
Ընդամենը 13-խաղաքարտանոց ձեռքեր՝ $\binom{52}{13} = 635{,}013{,}559{,}600$:

$\boxed{635{,}013{,}559{,}600 - 222{,}766{,}089{,}260 = 412{,}247{,}470{,}340}$

\subsubsection*{Խնդիր 4}
Բաժանել 10 մարդու 5 հոգանոց երկու թիմի:

Քանի որ երկու թիմերը չկարգավորված են (անզանազանելի պիտակներ).
$\dfrac{1}{2}\binom{10}{5} = \dfrac{252}{2}$:

$\boxed{126}$

\subsubsection*{Խնդիր 5}
$A=\{a,b,c,d,e,f,g,h\}$: 5-տարրանոց ենթաբազմությունների քանակը, որոնք պարունակում են $a$ տարրը:

Ֆիքսենք $a$-ն ենթաբազմության մեջ; ընտրենք ևս 4 տարր մնացած 7 տարրերից:

$\boxed{\binom{7}{4} = 35}$

\subsubsection*{Խնդիր 6}
$x^6 y^4$-ի գործակիցը $(2x+3y)^{10}$-ի վերլուծության մեջ:

Ըստ Նյուտոնի երկանդամի բանաձևի՝ $\binom{10}{6}(2x)^6(3y)^4 = \binom{10}{6}\cdot 2^6\cdot 3^4\cdot x^6 y^4$:

$\binom{10}{6}=210$, $2^6=64$, $3^4=81$:

$\boxed{210\cdot 64\cdot 81 = 1{,}088{,}640}$

\subsubsection*{Խնդիր 7}
Նավակը a1-ից հասնում է e5՝ շարժվելով միայն աջ կամ վերև:

a1-ից e5՝ 4 քայլ աջ ($a\to b\to c\to d\to e$) և 4 քայլ վերև ($1\to2\to3\to4\to5$): Ընդամենը 8 քայլ, ընտրել 4-ը աջ գնալու համար:

$\boxed{\binom{8}{4} = 70}$

\subsubsection*{Խնդիր 8}
Ցանցային ճանապարհներ $O(0,0)$-ից $A(7,8)$՝ անցնելով $(3,4)$-ով:

$O\to P(3,4)$: $\binom{7}{3}=35$ ճանապարհ: $P\to A$: $\binom{8}{4}=70$ ճանապարհ:

$\boxed{35\times 70 = 2450}$

\subsubsection*{Խնդիր 9}
6-ը երեք գումարելիների տրոհելու եղանակների քանակը:

\textbf{Կարգավորված} (կոմպոզիցիաներ, մասերը $\ge 1$). $\binom{5}{2}=\boxed{10}$:

\textbf{Չկարգավորված} (ամբողջ թվի տրոհումներ 3 մասի). $4+1+1$, $3+2+1$, $2+2+2$: Ընդամենը՝ $\boxed{3}$:

\subsubsection*{Խնդիր 10}
Գնել 7 թխվածք 3 տեսակից (զուգորդություններ կրկնություններով):

$\boxed{\binom{7+3-1}{3-1} = \binom{9}{2} = 36}$

\subsubsection*{Խնդիր 11}
``matematika'' բառի վերադասավորումները (10 տառ՝ a$\times 3$, m$\times 2$, t$\times 2$, e$\times 1$, i$\times 1$, k$\times 1$):

$\boxed{\frac{10!}{3!\cdot 2!\cdot 2!} = \frac{3{,}628{,}800}{24} = 151{,}200}$

\subsubsection*{Խնդիր 12}
Դասավորել 15 մաթեմատիկայի, 16 ինֆորմատիկայի, 12 ֆիզիկայի գրքեր դարակի վրա (նույն առարկայի գրքերը անզանազանելի են): Ընդամենը $15+16+12=43$ գիրք:

$\boxed{\frac{43!}{15!\cdot 16!\cdot 12!} = \binom{43}{15,\,16,\,12}}$

\subsubsection*{Խնդիր 13}
$x_1^4 x_2^3 x_3^2$-ի գործակիցը $(x_1+x_2+x_3)^9$-ի վերլուծության մեջ:

Ըստ բազմանդամի (multinomial) թեորեմի՝ $\dfrac{9!}{4!\,3!\,2!}$:

$\boxed{\frac{9!}{4!\,3!\,2!} = \frac{362880}{24\cdot 6\cdot 2} = 1260}$